% ============================================================
%  Information-theoretic message authentication
% ============================================================

\section{Information-theoretic message authentication}

\paragraph{From secrecy to integrity.} Shannon's theorem tells us that perfect secrecy is possible but expensive: the key must be at least as long as the message. Even if we accept this cost, the One-Time Pad gives us \textbf{confidentiality only}. It is malleable: by flipping bits of the ciphertext, Eve flips the same bits of the message, and Bob cannot notice:
\[
  \Dec(k, c \oplus \Delta) = m \oplus \Delta
\]
So, still \textbf{without any computational assumption}, the natural next question is: can we guarantee \textbf{integrity}, so that Bob always detects a modified message? This leads to \textbf{information-theoretic message authentication codes}. We will see that the story repeats itself: integrity against an unbounded adversary is possible, but the key must grow with the number of messages we want to authenticate.

\subsection{The setting}

In this section we \textbf{do not care about confidentiality}: the message travels in clear. What we want is \textbf{integrity}, that is, Bob must be able to detect whether the message was modified on the way.

Alice and Bob share a secret key $K$, chosen uniformly. Alice sends the message together with a \textbf{tag}, computed with a \textbf{keyed} function:
\[
  \tau = \Tag(K, m), \qquad \Tag : \Kspace \times \Mspace \to \Tspace
\]
where $\Tspace$ is the \textbf{tag space}. Here $\Tag$ is deterministic, and usually it is a keyed hash function, $\Tag(K, m) = h_K(m)$.

\begin{figure}[H]
  \centering
  \begin{tikzpicture}[box/.style={draw, rounded corners=2pt, minimum height=7mm, inner sep=4pt}]
    \node (A) at (0,0) {Alice};
    \node[box] (T) at (2.8,0) {$\Tag(K, m)$};
    \node[box] (V) at (8.2,0) {$\Tag(K, m) \overset{?}{=} \tau$};
    \node (B) at (11,0) {Bob};
    \node (K1) at (2.8,-1.2) {$K$};
    \node (K2) at (8.2,-1.2) {$K$};
    \draw[msg] (A) -- node[above] {$m$} (T);
    \draw[msg] (T) -- node[above] {$(m, \tau)$} (V);
    \draw[msg] (V) -- (B);
    \draw[msg] (K1) -- (T);
    \draw[msg] (K2) -- (V);
  \end{tikzpicture}
  \caption{An integrity-only MAC: the message travels in clear, together with its tag.}
\end{figure}

The protocol, step by step:
\begin{enumerate}
  \item \textbf{Alice} computes the tag of the message, $\tau = \Tag(K, m)$.
  \item Alice sends the \textbf{pair} $(m, \tau)$.
  \item \textbf{Bob} receives $(m, \tau)$, \textbf{recomputes} $\Tag(K, m)$ with the shared key and compares it with $\tau$: if they match he accepts the message, otherwise he rejects it.
\end{enumerate}

\paragraph{The threat.} Eve is an \textbf{active} adversary. She intercepts $(m, \tau)$ and tries to replace it with a pair $(m^*, \tau^*)$, with $m^* \neq m$, that Bob accepts. This is a \textbf{forgery}. The scheme is good if forging is hard \textbf{even after seeing valid pairs}.

\begin{framedobs}{Why the key is essential}
  Suppose we used a hash function \textbf{without a key}, a fixed public function $h$, and set $\tau = h(m)$. Eve intercepts $(m, \tau)$, replaces the message with $m^*$, computes $h(m^*)$ herself (the function is public) and sends $(m^*, h(m^*))$. Bob recomputes $h(m^*)$, sees that it matches, and accepts the forged message.

  With a key, instead, the key \textbf{selects which function} of the family $\{h_k\}$ is used. Eve knows the family (Kerckhoffs's principle) but not which function was chosen, so she cannot compute the correct tag for $m^*$.
\end{framedobs}

Since only the key holders can compute valid tags, a valid tag guarantees both \textbf{integrity} (the message was not modified) and \textbf{authenticity} (it was produced by someone holding the key, that is, by Alice or Bob). It is the shared secret key that gives both properties together.

\subsection{The security definition}

\begin{frameddefn}[label={def:mac-security}]{$t$-time $\varepsilon$-statistical security}
  A MAC $\Pi = (\Tag)$ is $t$-time $\varepsilon$-statistically secure if, for all \textbf{distinct} messages $m, m_1, \ldots, m_t \in \Mspace$ and all tags $\tau, \tau_1, \ldots, \tau_t \in \Tspace$:
  \[
    \Pr_K\big[\Tag(K, m) = \tau \;\big|\; \Tag(K, m_1) = \tau_1, \ldots, \Tag(K, m_t) = \tau_t\big] \leq \varepsilon
  \]
\end{frameddefn}

In words: Eve has seen $t$ valid pairs $(m_i, \tau_i)$. For any \textbf{new} message $m$ and any tag $\tau$ she proposes, the probability that $\tau$ is the correct tag is at most $\varepsilon$.
\begin{itemize}
  \item $t$ $\to$ how many valid pairs the adversary sees (bigger is better);
  \item $\varepsilon$ $\to$ the forging probability (smaller is better);
  \item the adversary is \textbf{unbounded}, which is why the security is called ``statistical''.
\end{itemize}

\begin{framedobs}[label={obs:eps-lower-bound}]{We can never get $\varepsilon = 0$}
  In fact $\varepsilon \geq \frac{1}{|\Tspace|}$ always. Fix everything except $\tau$. The tag of $m$ takes some value, so the probabilities over all possible tags sum to $1$:
  \[
    \sum_{\tau \in \Tspace} \Pr\big[\Tag(K, m) = \tau \mid \ldots\big] = 1
  \]
  A sum of $|\Tspace|$ terms equal to $1$ has at least one term $\geq \frac{1}{|\Tspace|}$, and guessing that tag succeeds with at least that probability. So, to get a small $\varepsilon$, we need long tags.
\end{framedobs}

We often write $\varepsilon = 2^{-\lambda}$, where $\lambda = \log_2 \frac{1}{\varepsilon}$ is the \textbf{security level in bits}: forging is as hard as guessing a random $\lambda$-bit string. For example, $\lambda = 128$ gives $\varepsilon = 2^{-128}$. With this notation, \Cref{obs:eps-lower-bound} says that tags must have \textbf{at least $\lambda$ bits}.

\begin{framedprob}[label={prob:otp-mac}]{The One-Time Pad is not a MAC}
  Let $\Mspace = \Kspace = \Tspace = \{0,1\}^n$ and $\Tag(k, m) = k \oplus m$. Show that this MAC is \textbf{not} $1$-time $\varepsilon$-statistically secure for any $\varepsilon < 1$.
\end{framedprob}

\begin{proof}[Solution]
  Think of the forgery game: Eve sees one valid pair $(m, \tau)$ with $\tau = \Tag(k, m)$, and must output $(m', \tau')$ with $m' \neq m$ and $\tau' = \Tag(k, m')$.
  \begin{enumerate}
    \item \textbf{Recover the key.} From $\tau = k \oplus m$, XOR-ing both sides with $m$ gives $k = \tau \oplus m$. One valid pair reveals the whole key.
    \item \textbf{Forge.} Choose any $m' \neq m$ and output $\tau' = (\tau \oplus m) \oplus m' = k \oplus m'$, which is exactly the correct tag of $m'$.
  \end{enumerate}
  Eve wins with probability $1$. In the language of \Cref{def:mac-security}: given $\Tag(K, m) = \tau$, the key is determined, so
  \[
    \Pr_K\big[\Tag(K, m') = \tau \oplus m \oplus m' \;\big|\; \Tag(K, m) = \tau\big] = 1
  \]
  and the scheme is $1$-time $\varepsilon$-secure only for $\varepsilon = 1$, that is, not secure at all.
\end{proof}

\textbf{The lesson.} The One-Time Pad is perfectly secret, yet it is useless for integrity. Secrecy asks that the output reveals nothing about the \textit{message}; a MAC needs more: after one tag, the tags of \textit{other} messages must still be unpredictable. Here the tag of $m$ and the tag of $m'$ are tied by the fixed relation $\Tag(k, m) \oplus \Tag(k, m') = m \oplus m'$. This is exactly what pairwise independence (below) rules out.

\subsection{The price: long keys (again)}

As with perfect secrecy, statistical security is expensive.

\begin{framedthm}[label={thm:mac-key-size}]{}
  Any $t$-time $2^{-\lambda}$-statistically secure MAC has a key of \textbf{at least} $(t+1)\lambda$ bits.
\end{framedthm}

The proof was omitted in class. The intuition: every valid pair Eve sees can ``use up'' about $\lambda$ bits of the key's uncertainty. To still have $\lambda$ bits of uncertainty about the \textbf{next} tag after $t$ pairs, the key needs $(t+1)\lambda$ bits, that is, $\lambda$ bits for every message to authenticate, plus $\lambda$. The key grows \textbf{linearly with the number of messages}: this is the MAC analogue of Shannon's theorem.

\subsection{Pairwise independent hash functions}

\begin{frameddefn}[label={def:pairwise}]{Pairwise independent hash family}
  A family of functions $\Hfam = \{h_k : \Mspace \to \Tspace\}_{k \in \Kspace}$ is \textbf{pairwise independent} if, for all $m \neq m'$, when $K$ is uniform the pair
  \[
    \big(h_K(m),\, h_K(m')\big) \text{ is uniform over } \Tspace \times \Tspace
  \]
  that is, $\Pr_K[h_K(m) = \tau \wedge h_K(m') = \tau'] = \frac{1}{|\Tspace|^2}$ for all $\tau, \tau'$.
\end{frameddefn}

\textbf{Intuition.} Pick a random function from the family. Its outputs on \textbf{any two different inputs} look like two independent, uniform values: knowing $h_K(m)$ tells you \textbf{nothing} about $h_K(m')$. The independence is between the outputs on two \textit{inputs}, not between two functions.

\textbf{Consequence: each output is uniform.} Summing the joint probability over $\tau'$:
\[
  \Pr_K[h_K(m) = \tau] = \sum_{\tau' \in \Tspace} \frac{1}{|\Tspace|^2} = \frac{1}{|\Tspace|}
\]

\begin{framedthm}[label={thm:pairwise-mac}]{}
  If $\Hfam$ is pairwise independent, then $\Tag(K, m) = h_K(m)$ is a \textbf{1-time $\frac{1}{|\Tspace|}$-statistically secure} MAC.
\end{framedthm}

\begin{proof}
  Eve sees one pair $(m, \tau)$ and tries $(m', \tau')$ with $m' \neq m$. By the definition of conditional probability:
  \[
    \Pr\big[\Tag(K, m') = \tau' \mid \Tag(K, m) = \tau\big]
    = \frac{\Pr[h_K(m') = \tau' \wedge h_K(m) = \tau]}{\Pr[h_K(m) = \tau]}
    = \frac{1/|\Tspace|^2}{1/|\Tspace|} = \frac{1}{|\Tspace|}
  \]
  The numerator uses pairwise independence (the pair is uniform over $\Tspace^2$); the denominator uses the consequence above (each output is uniform). Eve's chance is exactly that of a random guess, which is optimal by \Cref{obs:eps-lower-bound}.
\end{proof}

\subsection{A concrete construction: $am + b \bmod p$}

\begin{framedprop}[label={prop:affine}]{}
  Let $p$ be prime, $\Mspace = \Tspace = \Z_p$, $\Kspace = \Z_p^2$, and
  \[
    h_{(a,b)}(m) = a \cdot m + b \pmod p
  \]
  Then $\Hfam = \{h_{(a,b)}\}_{(a,b) \in \Z_p^2}$ is pairwise independent.
\end{framedprop}

\begin{proof}
  Fix $m \neq m'$ and $\tau, \tau'$. We count how many keys give tag $\tau$ to $m$ \textbf{and} tag $\tau'$ to $m'$. The two conditions form a linear system in the unknowns $a, b$:
  \[
    \begin{pmatrix} m & 1 \\ m' & 1 \end{pmatrix} \begin{pmatrix} a \\ b \end{pmatrix} = \begin{pmatrix} \tau \\ \tau' \end{pmatrix}
  \]
  The determinant is $m - m' \neq 0$. Since $p$ is \textbf{prime}, $\Z_p$ is a field, so every nonzero element is invertible and the matrix is invertible. The system has \textbf{exactly one} solution $(a, b)$, and the key is uniform over $p^2$ possibilities:
  \[
    \Pr_{(a,b)}\big[h(m) = \tau \wedge h(m') = \tau'\big] = \frac{1}{p^2} = \frac{1}{|\Tspace|^2}
  \]
  It is the same ``exactly one key'' argument as for the One-Time Pad.
\end{proof}

\paragraph{Why $p$ must be prime.} To solve the system we must \textbf{divide} by $m - m'$. Modulo a prime, every nonzero number has an inverse. Modulo a composite number this fails: modulo $6$, for example, $2$ has no inverse, and the system could have zero or several solutions.

\paragraph{Geometric intuition.} $y = ax + b$ is the equation of a \textbf{line}. Choosing the key means choosing a \textbf{random line} over $\Z_p$, and the tag of $m$ is the \textbf{height of the line at the point $m$}. Every message-tag pair that Eve sees reveals \textbf{one point} of the line.

\begin{framedobs}{The key can be used only once}
  As with the One-Time Pad, this MAC is secure only if each key authenticates \textbf{a single message}.
  \begin{itemize}
    \item \textbf{One pair} $\to$ one point. Exactly $p$ lines pass through it, one for each slope, and each gives a different tag for any other message $m^*$. So Eve can only guess, with probability $\frac{1}{p}$.
    \item \textbf{Two pairs} $\to$ two points, and only \textbf{one line} passes through two points. Eve solves the system, recovers the key, and can forge a valid tag for any message.
  \end{itemize}
  The key has $2\log p$ bits, that is $(t+1)\lambda$ with $t = 1$ and $\lambda = \log p$: by \Cref{thm:mac-key-size}, it is exactly enough for one message, and no more.
\end{framedobs}

\begin{framedex}{An attack with two pairs}
  With $p = 7$ and secret key $(a, b) = (3, 5)$, Eve sees:
  \begin{itemize}
    \item $m_1 = 2$, with $\tau_1 = 3 \cdot 2 + 5 = 11 \equiv 4$;
    \item $m_2 = 4$, with $\tau_2 = 3 \cdot 4 + 5 = 17 \equiv 3$.
  \end{itemize}
  Subtracting the two equations removes $b$:
  \[
    a = \frac{\tau_1 - \tau_2}{m_1 - m_2} = \frac{4 - 3}{2 - 4} = \frac{1}{-2} \equiv 1 \cdot 5^{-1} \equiv 3 \pmod 7,
    \qquad b = \tau_1 - a m_1 = 4 - 6 \equiv 5 \pmod 7
  \]
  Here $-2 \equiv 5$ and $5^{-1} \equiv 3$ because $5 \cdot 3 = 15 \equiv 1 \pmod 7$. Eve has recovered the whole key $(3, 5)$.

  With a single pair, instead, Eve only knows that $2a + b \equiv 4$: this equation has $7$ solutions, one for each value of $a$, and each gives a different tag for a new message.
\end{framedex}

\subsection{Beyond the lecture: $t$ messages with polynomials}

To authenticate more messages, replace the line with a richer curve: a random \textbf{polynomial of degree $t$},
\[
  h_{(a_0, \ldots, a_t)}(m) = a_0 + a_1 m + a_2 m^2 + \cdots + a_t m^t \pmod p
\]
where the key is the list of the $t+1$ coefficients. The geometry generalizes: through $2$ points passes a single line, through $3$ points a single parabola, and through $t+1$ points a single polynomial of degree $t$. The reason is the Vandermonde matrix.

\begin{frameddefn}{Vandermonde matrix}
  The Vandermonde matrix of the points $x_0, \ldots, x_t$ is the matrix whose $i$-th row contains the powers of $x_i$:
  \[
    V = \begin{pmatrix}
      1 & x_0 & x_0^2 & \cdots & x_0^t \\
      1 & x_1 & x_1^2 & \cdots & x_1^t \\
      \vdots & \vdots & \vdots & & \vdots \\
      1 & x_t & x_t^2 & \cdots & x_t^t
    \end{pmatrix},
    \qquad \det V = \prod_{i < j} (x_j - x_i)
  \]
  The determinant is nonzero if and only if the points are distinct. Hence through $t+1$ distinct points passes \textbf{exactly one} polynomial of degree $t$: its coefficients are the unique solution of $V \cdot (a_0, \ldots, a_t)^\top = (\tau_0, \ldots, \tau_t)^\top$.
\end{frameddefn}

For $t = 1$ the Vandermonde matrix is $\begin{pmatrix} 1 & m \\ 1 & m' \end{pmatrix}$: it is the matrix of \Cref{prop:affine} with the columns swapped, which only changes the sign of the determinant.

Consequences:
\begin{itemize}
  \item after seeing $t$ pairs, $p$ polynomials are still possible and the next tag is uniform $\to$ this is a \textbf{$t$-time $\frac{1}{p}$-secure MAC};
  \item with $t+1$ pairs Eve reconstructs the polynomial, that is, the key;
  \item the key has $(t+1)\log p$ bits, which matches the lower bound $(t+1)\lambda$ of \Cref{thm:mac-key-size} with $\lambda = \log p$: the construction is \textbf{optimal}.
\end{itemize}
